\documentclass[uplatex,dvipdfmx]{u-sho_jsarticle}
\usepackage{u-sho_style}
\usepackage{u-sho_math}
\hypersetup{
  pdfinfo={
    CreationDate={D:20201231210407+09'00'},
    ModDate={\pdffilemoddate{\jobname}},
    Subject={upLaTeX template with u-sho's packages},
    Keywords={upLaTeX template;u-sho jsarticle;u-sho math;u-sho style}
  }
}

\title{\upLaTeXe テンプレート} % chktex 1
\author{u-sho (上地 将平)}
\date{作成：令和2年12月31日\hspace{2zw}更新：\和暦\today}%chktex 19

\begin{document}
  \maketitle
  \tableofcontents

  \clearpage

  \part{ho}
    \section{目的}
      知らんが

    \section{概要}
      参考文献~\cite{キー1}

  \part{he}
    \section{鳴き声} %chktex 19

      \subsection{にゃあにゃあ}
        なんか

      \subsubsection{\texorpdfstring{\bgOrange{\tWhite{がお}}がお}{がおがお}}
        すると

      \subsection{ぴえん} % chktex 19

      \subsubsection{ぱおーん}
        \fcolorbox{red}[HTML]{888888}{\tWhite{ぱおーん}}
        \prettyref{th:center}や\prettyref{tb:result}，\prettyref{fg:paopao}，\prettyref{eq:newpage}を \verb|\prettyref| コマンドで引いてみた

  \newpage

    \section{結果および考察}
      \subsection{がおがお}
        しらんけど\num{6e2}回サイコロふった．初期位相\ang{-66;5;4}だったとかなんとか

        \begin{ctable}[h]
          \caption{サイコロを振った結果（\(N=600\)）}\label{tb:result}
          \begin{tabular}{%
            c|S[table-format=3.0]cS[table-format=+1.1e+1,table-figures-uncertainty=1,color=orange]%
          }
            \hline
            出た目 & 回数\(n\) & 割合\(n/N\) & {\color{orange}\((n/N - 1/6)\)} \\
            \hline
            1      &  88       & 0.147       & -2.0+-1  e-2 \\
            2      & 102       & 0.170       &  3  +-0.1e-3 \\
            \cline{2-4}
            3      &  84       & 0.140       & -2.7     e-2 \\
            4      & 114       & 0.190       &  2.3     e-2 \\
            5      & 109       & 0.182       &  1.5     e-2 \\
            6      & 103       & 0.172       & \color{red}5e-3 \\
            \hline
          \end{tabular}
        \end{ctable}

      \subsection{ぱおぱお}
        これはなんの\emph{ず}？pyo
          \begin{figure}[!ht]
            \centering
              \includegraphics{img/350x150.png} %chktex 29
              \caption{ろんりかいろ}\label{fg:paopao}
          \end{figure}
        \numlist{1;2;10}番目のわちゃわちゃ\\
        \(1,\,2,\,\ldots\,,\,10\)あるいは\(1,\,2\dotsc 10\)．\(\llistto{Y}{m}\)は \(\llistto{X}{n}\)

  \newpage

        \subsubsection{にゅ〜ぺーじ}
          ひるべるとくうかん \(\lrp[\big]{\H, \L\H} \lrb[\Big]{x,y}\):
          \begin{equation}\label{eq:newpage}
            \sfrac{3}{100}\times \mathcal{X}\inp*{\wfrac{100}{3}}{3+i} %\hsp*{3+i}{\frac{4}{101}}
            \Exist{ρ},\All{\rho}, \vphantom{\rho}^\forall\!\rho \bigl\langle\!\!\bigl\langle
          \end{equation}
          あるいは
          \begin{theorem}[中間値の定理]\label{th:center}
            閉区間\(\closedInterval{a,b}\subset\R\)について連続な関数\(f:\R\to\R\)を考える．
            \(f\lrp{a}\leq f\lrp{b}\)であるとき，
            \(\openInterval{f\lrp{a},f\lrp{b}}\)上の任意の\(k\in\R\)に対して，
            \(f\lrp{c}=k\)なる\(c\in\openInterval{a,b}\)が存在し，
            \(f\lrp{a}>f\lrp{b}\)であるとき，
            \(\openInterval{f\lrp{b},f\lrp{a}}\)上の任意の\(k\in\R\)に対して，
            \(f\lrp{c}=k\)なる\(c\in\openInterval{b,a}\)が存在する．
          \end{theorem}

  \clearpage

  \begin{thebibliography}{9}
    \bibitem{キー1} 参考文献の名前・著者1 available at \url{https://github.com/u-sho/upLaTeX-template}
  \end{thebibliography}

\end{document}
